\chapter{Conclusion and policy}\label{ch:concl}

\section{What binds}
\begin{table}[h]\centering\small
\caption{The pre-specified tests}\label{tab:summary}
\begin{tabularx}{\textwidth}{>{\raggedright\arraybackslash}p{4.2cm}>{\raggedright\arraybackslash}p{3.2cm}L}
\toprule
Hypothesis & Pre-specified result & Comment \\
\midrule
\multicolumn{3}{l}{\emph{Primary family (Holm correction at 5\%)}} \\
H8: trading leaves futures in a ban & Supported ($p<10^{-10}$) & Identity holds under intraday enforcement \\
H3: costs set the band edges & Supported ($p<10^{-10}$) & Placebo-calibrated evidence weak; mostly a common shift \\
H6: short-sale access & Supported ($p=5\times10^{-6}$) & Through price delay; drift goes the other way \\
H7: one-sided correction & Rejected, opposite sign & Survives a noise correction \\
\midrule
\multicolumn{3}{l}{\emph{Secondary (unadjusted)}} \\
Level effect by composition & $p=0.083$ & Signs as predicted; median split $p=0.025$ \\
Density at 95\% & No discontinuity & No bunching \\
Regression discontinuity & Basis $-12.9$\bp, $p=0.02$ & Gap present before the ban \\
Hysteresis & $+0.4$\bp, $p=0.64$ & No extra level shift \\
H9 anticipation & $p=0.019$ & Upper tail grows toward 95\% \\
H9 reform & $p=0.11$ & Right direction \\
H2 bad news in bans & $p=0.59$ & No power \\
H1 frictions and band & $p=0.002$ & Across stocks only \\
H4 capital & $p=0.11$ & Clearer in differences \\
H5 anomalies & $p=0.57$ & Not supported value-weighted \\
\bottomrule
\end{tabularx}
\end{table}

The answer to the question this thesis asks is uneven across the limits. Position limits bind hard on trading: the ban
moves trading out of futures as soon as it is enforced on intraday positions. Their effect on prices is real but not the
one the reduce-only arbitrage channel predicts; the evidence points instead to the blocking of new demand at the start of
a ban. Short-sale access matters for how fast prices adjust to market news, both when a stock enters the F\&O segment and
across all stocks, but less clearly for overpricing. Trading costs shape the band, which is wide and asymmetric, but its
edges do not move one-for-one with the tax: the full cost of a new position lies in the outer tail of the basis, and the
cost of keeping arbitrage capital in place is small. Arbitrage capital leaves weak traces in the predicted direction.

With ten secondary statistics reported without adjustment, one or two values below 0.05 would be expected by chance. The
secondary results should be read as a pattern, not one by one.

\section{Policy}
Position limits and the ban exist to contain concentrated positions and manipulation in a market where speculative
demand is large. SEBI's study of individual traders in equity derivatives (SEBI, 2024a) shows how large: over the three
years to March 2024, 93\% of more than one crore individual traders lost money, on average about Rs 2 lakh each including
transaction costs, and their aggregate losses exceeded Rs 1.8 lakh crore. In the last of those years proprietary traders and
foreign portfolio investors earned gross profits, 96 to 97\% of them from algorithmic trading. SEBI's measures of October 2024 for index
derivatives cite the rise in retail participation and speculative volumes (SEBI, 2024b), and its measures of May 2025
changed how open interest is measured and how position limits are set (SEBI, 2025). Most of that
activity is in index options, not in the single-stock futures studied here, but the same pressure of new speculative
demand is what pushes a stock towards its limit.

The results bear on the trade-off in four ways. First, the ban is costly for liquidity: futures trading falls by almost
90\% in a ban, while cash trading rises a little. Second, its main price effect appears at the start of a ban, where it
fits the blocking of new demand, and fades within about a week; once a stock is in the ban, staying in it does not move
the basis further. If that reading is right, a ban restrains the demand it is designed to restrain, but at the cost of the
arbitrage that would otherwise keep futures and cash prices together. Third, enforcement matters: the rule changes trading in the way it
is meant to only when positions are checked during the day. The October 2025 rule, which allows any trade that lowers an
entity's end-of-day delta, has visibly loosened the link between volume and open interest in banned stocks.

Fourth, the 2025 reform met its stated aim of fewer ``spurious'' bans (SEBI, 2025, para.\ 4.2). Before the reform, bans
were not caused by falls in the limit: almost every entry came from rising open interest. What made many of them avoidable
was the measure of open interest, which counted every option at its full notional value. Measuring positions in delta terms
removes almost all of these bans, the new and lower limits add some back, and the combination cut ban days by about three
quarters in the first year. Whether the new formula has changed the cost of each ban needs more post-reform data than the
44 spells that began after the reform. The 2020 restoration of the limit adds one more lesson: a looser limit lowered
utilisation and bans and made a rich basis less common, so the level of the limit affects prices even before a ban starts.

The SLB market remains thin. Because only F\&O stocks may be sold short and the overnight short is in practice a stock
future, entry into and exit from the F\&O segment change the short-sale constraint of a stock at once. The faster price
adjustment after entry is one measurable benefit of that access.

\section{Limitations and further work}
The data are daily, so intraday arbitrage and the exact timing of trades within the day are outside the scope, and noise
in daily prices inflates measured reversion. Arbitrage inventory is not observed. The post-reform period is short. Three
papers relevant to positioning (Saurav, Srivastava and Agarwalla, 2026; Gagnon, 2018; Brennan and Schwartz, 1990) could
not be read in full. Intraday data would allow sharper tests of the ban's effect on correction, and a longer post-reform
sample would allow a proper test of the 2025 formula.
